\par
\clearpage
\SourceLeaf{129}\GaussVTheoryEarlyFolio{125}
\noindent überaus beschwerlichen Rechnungen für eine bedeutend grössere Zahl der Ausführbarkeit unübersteigliche Schwierigkeiten entgegensetzen würden. Das Beste wird sein, diesen Weg ganz zu verlassen, der unwillkürlich an die Versuche erinnert. die Planetenbewegungen durch immer mehr gehäufte Epicykeln zu erklären.

In der gegenwärtigen Abhandlung werde ich die allgemeine Theorie des Erdmagnetismus, unabhängig von allen besondern Hypothesen über die Vertheilung der magnetischen Flüssigkeit im Erdkörper, entwickeln, und zugleich die Resultate mittheilen, welche ich aus der ersten Anwendung der Methode erhalten habe. So unvollkommen diese Resultate auch sein müssen, so werden sie doch einen Begriff davon geben können, was man hoffen darf in Zukunft zu erreichen, wenn einer feinern und wiederholten Ausfeilung derselben erst zuverlässige und vollständige Beobachtungen aus allen Gegenden der Erde werden untergelegt werden können.

\begin{center}\rule{.22\linewidth}{.35pt}\end{center}
\Article{1}
Die Kraft, welche einer in ihrem Schwerpunkte aufgehängten Magnetnadel an jedem Orte der Erde eine bestimmte Richtung ertheilt, indem jede fremde äussere Ursache, die auf die Nadel wirken könnte (wie die Nähe eines andern künstlichen Magnets, oder die Nähe des Leiters eines galvanischen Stroms) als beseitigt vorausgesetzt wird, nennt man die erdmagnetische Kraft, insofern man den Sitz ihrer Ursache nur in dem Erdkörper selbst suchen kann. Zweifelhaft ist allerdings, ob die regelmässigen und unregelmässigen stündlichen Aenderungen in jener Kraft nicht ihre nächsten Ursachen ausserhalb des Erdkörpers haben mögen, und es steht zu hoffen, dass die jetzt auf diese Erscheinungen allgemein gerichtete Aufmerksamkeit der Naturforscher uns darüber in Zukunft bedeutende Aufschlüsse geben werde. Allein man darf nicht vergessen, dass diese Änderungen vergleichungsweise nur sehr klein sind, und dass also eine viel stärkere beharrlich wirkende Hauptkraft da sein muss, deren Sitz wir in der Erde selbst annehmen. Es knüpft sich hieran sofort die Folgerung, dass die zur Untersuchung dieser Hauptkraft dienenden thatsächlichen Grundlagen eigentlich von den erwähnten anomalischen Aenderungen befreiet sein sollten, was nur durch Mit-\par
\clearpage
\SourceLeaf{130}\GaussVTheoryEarlyFolio{126}
\noindent telwerthe aus zahlreichen fortgesetzten Beobachtungen möglich ist, und dass so lange solche reine Resultate nicht von einer grossen Anzahl von Punkten auf der ganzen Erdoberfläche vorhanden sind, das Höchste, was man wird erreichen können, eine Annäherung ist, wobei Differenzen von der Ordnung solcher Anomalien zurückbleiben können.

\Article{2}
Die Grundlage unserer Untersuchungen ist die Voraussetzung, dass die erdmagnetische Kraft die Gesammtwirkung der magnetisirten Theile des Erdkörpers ist. Das Magnetisirtsein stellen wir uns als eine Scheidung der magnetischen Flüssigkeit vor: diese Vorstellungsweise einmal angenommen, gehört die Wirkungsart dieser Flüssigkeiten (Abstossung oder Anziehung des Gleichnamigen oder Ungleichnamigen im verkehrten Verhältniss des Quadrats der Entfernung) zu den erwiesenen physikalischen Wahrheiten. Eine Vertauschung dieser Vorstellungsart mit der \textsc{Ampère}schen, wonach, mit Beseitigung der magnetischen Flüssigkeiten, der Magnetismus nur in beharrlichen galvanischen Strömungen in den kleinsten Theilen der Körper besteht, würde in den Resultaten gar nichts abändern; dasselbe würde auch gelten, wenn man den Erdmagnetismus einer gemischten Ursache zuschreiben wollte, so dass derselbe theils aus Scheidung der magnetischen Flüssigkeiten in der Erde, theils aus galvanischen Strömungen in derselben herrührte, indem bekanntlich anstatt eines jeden galvanischen Stromes eine solche bestimmte Vertheilung von magnetischen Flüssigkeiten an einer von der Stromlinie begrenzten Fläche substituirt werden kann, dass dadurch in jedem Punkte des äussern Raumes genau dieselbe magnetische Wirkung ausgeübt wird, wie durch den galvanischen Strom.

\Article{3}
Zur Abmessung der magnetischen Flüssigkeiten legen wir, wie in der Schrift \textit{Intensitas vis magneticae} etc., diejenige Quantität nordlichen Fluidums als positive Einheit zum Grunde, welche auf eine eben so grosse Quantität desselben Fluidums in der zur Einheit angenommenen Entfernung eine bewegende Kraft ausübt, die der zur Einheit angenommenen gleich ist. Wenn wir von der magnetischen Kraft, welche in irgend einem Punkte des Raumes, als Wirkung von anderswo befindlichem magnetischen Fluidum, schlechthin sprechen, so ist darunter immer \par
\clearpage
\SourceLeaf{131}\GaussVTheoryEarlyFolio{127}
\noindent die bewegende Kraft verstanden, welche daselbst auf die Einheit des positiven magnetischen Fluidums ausgeübt wird. In diesem Sinne übt folglich die in Einem Punkt concentrirt gedachte magnetische Flüssigkeit $\mu$ in der Entfernung $\rho$ die magnetische Kraft $\frac{\mu}{\rho\rho}$ aus, und zwar abstossend oder anziehend in der Richtung der geraden Linie $\rho$, je nachdem $\mu$ positiv oder negativ ist. Bezeichnet man durch $a,b,c$ die Coordinaten von $\mu$ in Beziehung auf drei unter rechten Winkeln einander schneidende Axen; durch $x,y,z$ die Coordinaten des Punkts, wo die Kraft ausgeübt wird, so dass $\rho=\BareRoot{((x-a)^2+(y-b)^2+(z-c)^2)}$, und zerlegt die Kraft den Coordinatenaxen parallel, so sind die Componenten
\[
\frac{\mu(x-a)}{\rho^3},\qquad \frac{\mu(y-b)}{\rho^3},\qquad \frac{\mu(z-c)}{\rho^3}
\]
welche, wie man leicht sieht, den partiellen Differentialquotienten von $-\frac{\mu}{\rho}$ nach $x,y$ und $z$ gleich sind.

Wirken ausser $\mu$ noch andere Theile magnetischen Fluidums, $\mu',\mu'',\mu'''$ u. s. w., concentrirt in Punkten, deren Entfernung von dem Wirkungsorte bezugweise $\rho',\rho'',\rho'''$ u. s. w. ist, so sind die Componenten der ganzen daraus resultirenden magnetischen Kraft, parallel mit den Coordinatenaxen, gleich den partiellen Differentialquotienten von
\[
-\left(\frac{\mu}{\rho}+\frac{\mu'}{\rho'}+\frac{\mu''}{\rho''}+\frac{\mu'''}{\rho'''}+\text{u. s. w.}\right)
\]
nach $x,y$ und $z$.

\Article{4}
Man übersieht hienach leicht, welche magnetische Kraft in jedem Punkte des Raumes von der Erde ausgeübt werde, wie auch die magnetischen Flüssigkeiten in derselben vertheilt sein mögen. Man denke sich das ganze Volumen der Erde, so weit es freien Magnetismus, d. i. geschiedene magnetische Flüssigkeiten enthält, in unendlich kleine Elemente zerlegt, bezeichne unbestimmt die in jedem Elemente enthaltene Menge freien magnetischen Fluidums mit $\mathrm{d}\mu$, wobei südliches stets als negativ betrachtet wird; ferner mit $\rho$ die Entfernung des $\mathrm{d}\mu$ von einem unbestimmten Punkte des Raumes, dessen rechtwinklige Coordinaten $x,y,z$ sein mögen, endlich mit $V$ das Aggregat der $\frac{\mathrm{d}\mu}{\rho}$ mit verkehrtem Zeichen durch die Gesammtheit aller magnetischen Theilchen der Erde erstreckt. oder es sei
\par
\clearpage
